% =============================================================================
% Sorotan dan Highlights
% - 3--5 poin untuk setiap bahasa.
% - Maksimum 100 karakter per poin, termasuk spasi dan tanda baca.
% - Tidak merujuk pustaka, gambar, atau tabel.
% =============================================================================

\ipbprelimheading{SOROTAN}

\noindent
\MakeUppercase{\NamaPenulis}. \JudulSkripsi\unskip. Dibimbing oleh
\MakeUppercase{\NamaPembimbingSatu} dan \MakeUppercase{\NamaPembimbingDua}

\begin{enumerate}[leftmargin=0.5cm,labelwidth=0.25cm,labelsep=0.25cm,
  align=left,itemsep=0pt,parsep=0pt,topsep=12pt,partopsep=0pt]
  \item Penelitian ini membahas [masukkan topik utama], dengan tujuan untuk
    [nyatakan tujuan secara jelas dan ringkas].
  \item Penelitian ini menggunakan [sebutkan metode] untuk
    menganalisis/menginvestigasi [sebutkan objek kajian].
  \item Hasil penelitian menunjukkan bahwa [ringkas temuan utama dengan tepat].
  \item Temuan ini memberikan kontribusi terhadap [bidang/implikasi], serta
    menawarkan wawasan baru mengenai [aspek spesifik].
  \item Tugas akhir ini menunjukkan pentingnya [tekankan signifikansi], dan
    menyarankan potensi penerapan pada [konteks relevan].
\end{enumerate}

\vspace{10pt}

\noindent
Sorotan/\textit{highlights} ditulis dalam satu spasi, disusun 3--5 poin, dan
tidak lebih dari satu halaman. Setiap poin ditulis maksimal 100 karakter
termasuk spasi dan tanda baca. Tidak diperbolehkan mengacu pustaka, gambar,
dan tabel. Singkatan hanya dikenalkan jika masih digunakan lagi dalam bagian
lain sorotan/\textit{highlights}.

\begin{otherlanguage}{english}
\itshape
\vspace{20pt}
\ipbprelimsubheading{HIGHLIGHTS}

\noindent
\MakeUppercase{\NamaPenulisInggris}. \JudulSkripsiInggris\unskip. Supervised by
\MakeUppercase{\NamaPembimbingSatu} and \MakeUppercase{\NamaPembimbingDua}

\begin{enumerate}[leftmargin=0.5cm,labelwidth=0.25cm,labelsep=0.25cm,
  align=left,itemsep=0pt,parsep=0pt,topsep=7pt,partopsep=0pt]
  \item This study addresses [insert key topic], aiming to [state objective
    clearly and concisely].
  \item The research employs [mention methodology] to analyze/investigate
    [mention subject of study].
  \item Results indicate that [summarize key findings with precision].
  \item The findings contribute to [mention the field/implications], offering
    new insights into [specific aspect].
  \item This final assignment demonstrates the importance of [highlight
    significance], suggesting potential applications in [relevant context].
\end{enumerate}
\end{otherlanguage}
